\documentclass[12pt, a4paper]{report}

% ─────────────────────────────────────────────────────────────────────────────
% PAQUETES
% ─────────────────────────────────────────────────────────────────────────────
\usepackage[utf8]{inputenc}
\usepackage[T1]{fontenc}
\usepackage[spanish]{babel}
\usepackage{geometry}
\usepackage{graphicx}
\usepackage{float}
\usepackage{hyperref}
\usepackage{xcolor}
\usepackage{listings}
\usepackage{booktabs}
\usepackage{longtable}
\usepackage{array}
\usepackage{enumitem}
\usepackage{fancyhdr}
\usepackage{tcolorbox}
\usepackage{amsmath}
\usepackage{caption}
\usepackage{subcaption}
\usepackage{multirow}
\usepackage{lmodern}
\usepackage{microtype}
\usepackage{mdframed}

% ─────────────────────────────────────────────────────────────────────────────
% CONFIGURACION DE PAGINA
% ─────────────────────────────────────────────────────────────────────────────
\geometry{top=2.5cm, bottom=2.5cm, left=3cm, right=2.5cm, headheight=14pt}

% ─────────────────────────────────────────────────────────────────────────────
% COLORES
% ─────────────────────────────────────────────────────────────────────────────
\definecolor{satriBlue}{RGB}{30, 60, 114}
\definecolor{satriCyan}{RGB}{0, 180, 216}
\definecolor{satriGray}{RGB}{100, 110, 130}
\definecolor{satriGreen}{RGB}{0, 180, 80}
\definecolor{satriRed}{RGB}{220, 50, 50}
\definecolor{codeBg}{RGB}{248, 249, 250}
\definecolor{codeBorder}{RGB}{200, 210, 220}

% ─────────────────────────────────────────────────────────────────────────────
% ESTILO DE CODIGO
% ─────────────────────────────────────────────────────────────────────────────
\lstdefinestyle{satricode}{
    backgroundcolor=\color{codeBg},
    basicstyle=\ttfamily\small,
    breaklines=true,
    captionpos=b,
    commentstyle=\color{satriGray}\itshape,
    frame=single,
    framerule=0.5pt,
    rulecolor=\color{codeBorder},
    keywordstyle=\color{satriBlue}\bfseries,
    stringstyle=\color{satriGreen},
    numbers=left,
    numberstyle=\tiny\color{satriGray},
    numbersep=8pt,
    showspaces=false,
    showstringspaces=false,
    tabsize=4,
    xleftmargin=15pt
}
\lstset{style=satricode}

% ─────────────────────────────────────────────────────────────────────────────
% ENCABEZADO Y PIE DE PAGINA
% ─────────────────────────────────────────────────────────────────────────────
\pagestyle{fancy}
\fancyhf{}
\fancyhead[L]{\small\color{satriGray}SATRI --- Sistema de Analisis y Triage de Incidentes}
\fancyhead[R]{\small\color{satriGray}\leftmark}
\fancyfoot[C]{\small\thepage}
\renewcommand{\headrulewidth}{0.4pt}

% ─────────────────────────────────────────────────────────────────────────────
% CAJAS INFORMATIVAS
% ─────────────────────────────────────────────────────────────────────────────
\tcbuselibrary{skins, breakable}
\newtcolorbox{notabox}[1][]{colback=satriBlue!5!white, colframe=satriBlue!60!white,
    fonttitle=\bfseries, title=Nota, #1}
\newtcolorbox{alertbox}[1][]{colback=satriRed!5!white, colframe=satriRed!60!white,
    fonttitle=\bfseries, title=Importante, #1}

% ─────────────────────────────────────────────────────────────────────────────
% COMANDO PARA ESPACIO DE IMAGEN (placeholder)
% ─────────────────────────────────────────────────────────────────────────────
\newcommand{\imagespace}[2]{%
\begin{figure}[htbp]
    \centering
    \fbox{%
        \begin{minipage}{0.85\textwidth}
            \centering
            \vspace{3cm}
            {\large\color{satriGray} [INSERTAR IMAGEN: #1]}\\[0.5em]
            {\small\color{satriGray}(Reemplazar este recuadro con la captura de pantalla correspondiente)}
            \vspace{3cm}
        \end{minipage}%
    }
    \caption{#2}
\end{figure}%
}

% ─────────────────────────────────────────────────────────────────────────────
% HYPERREF
% ─────────────────────────────────────────────────────────────────────────────
\hypersetup{
    colorlinks=true,
    linkcolor=satriBlue,
    urlcolor=satriCyan,
    citecolor=satriBlue,
    pdftitle={Documentacion Tecnica SATRI},
    pdfauthor={Equipo SATRI},
    pdfsubject={Ciberseguridad - Arquitectura de Microservicios}
}

% ─────────────────────────────────────────────────────────────────────────────
% CONTROL DE FLOATS (evita paginas en blanco por figuras huerfanas)
% ─────────────────────────────────────────────────────────────────────────────
\renewcommand{\topfraction}{0.9}
\renewcommand{\bottomfraction}{0.8}
\renewcommand{\textfraction}{0.05}
\renewcommand{\floatpagefraction}{0.7}
\setcounter{totalnumber}{6}

% =============================================================================
% INICIO DEL DOCUMENTO
% =============================================================================
\begin{document}
\chapter{Componentes de Telemetria (Endpoints)}
% =============================================================================

\section{Agente Desktop SATRI}

\subsection{Descripcion}
El Agente Desktop es una aplicacion \textbf{Python compilada con PyInstaller} que se ejecuta como proceso en segundo plano en equipos Windows. El resultado de la compilacion es un ejecutable monolitico \texttt{SATRI\_Protection.exe} que los usuarios finales instalan, sin necesitar Python en su equipo.

\subsection{Modulos de Monitoreo}

\begin{table}[H]
\centering
\caption{Modulos del Agente Desktop SATRI y sus funciones}
\begin{tabular}{@{}p{3.5cm}p{4cm}p{5cm}@{}}
\toprule
\textbf{Modulo} & \textbf{Archivo} & \textbf{Funcion} \\
\midrule
Process Monitor     & \texttt{process\_monitor.py}   & Detecta procesos maliciosos y no autorizados \\
Ransomware Guard    & \texttt{ransomware\_guard.py}  & Monitorea cifrado masivo de archivos (FIM) \\
Download Scanner    & \texttt{download\_scanner.py}  & Escanea descargas por extensiones sospechosas \\
WiFi Check          & \texttt{wifi\_check.py}        & Detecta redes WiFi inseguras y portales cautivos \\
Network Monitor     & \texttt{network\_monitor.py}   & Monitorea conexiones de red salientes \\
Activity Monitor    & \texttt{activity\_monitor.py}  & Captura metricas de comportamiento para UEBA \\
Browser Bridge      & \texttt{browser\_bridge.py}    & Recibe eventos de la extension de navegador \\
Kernel Bridge       & \texttt{kernel\_bridge.py}     & Interfaz con el driver Ring-0 (opcional) \\
IOC Sync            & \texttt{ioc\_sync.py}          & Sincroniza lista de IoCs locales desde el backend \\
Reporter            & \texttt{reporter.py}           & Envia eventos al MML via REST con autenticacion \\
Local API           & \texttt{local\_api.py}         & API local en puerto 18001 para comandos SOAR \\
\bottomrule
\end{tabular}
\end{table}

\subsection{Flujo de Inicio y Consentimiento}
Antes de comenzar cualquier monitorizacion, el agente verifica si existe un registro de consentimiento del usuario en disco. Si no existe, muestra la pantalla de \textit{onboarding} con el texto legal completo (descrito en el Capitulo \ref{privacidad}) y espera la aceptacion explicita antes de arrancar los modulos de vigilancia.

\begin{lstlisting}[language=Python, caption={Flujo de inicio del Agente Desktop (main.py)}]
if not verify_user_consent():
    logger.error("Consentimiento no otorgado. Saliendo...")
    exit(1)

# Iniciar modulos como threads daemon
threading.Thread(target=start_monitor, daemon=True).start()
threading.Thread(target=start_ransomware_guard, daemon=True).start()
threading.Thread(target=start_scanner, daemon=True).start()
threading.Thread(target=start_wifi_check, daemon=True).start()
threading.Thread(target=start_network_monitor, daemon=True).start()
threading.Thread(target=start_activity_monitor, daemon=True).start()
threading.Thread(target=start_browser_bridge, daemon=True).start()

# El hilo principal corre la API local (para comandos de aislamiento)
run_local_api()
\end{lstlisting}

\imagespace{SATRI Protection.exe corriendo en el Administrador de Tareas de Windows}{Agente Desktop SATRI: ejecucion en segundo plano en Windows}

\section{Driver de Kernel Windows (Ring-0)}

\subsection{Descripcion}
SATRI incluye un \textbf{driver de kernel para Windows} implementado en C usando el \textit{Windows Driver Model} (WDM). Al operar en el nivel Ring-0, puede interceptar llamadas al sistema operativo que los procesos de usuario (Ring-3) no pueden monitorear, lo que permite detectar rootkits y amenazas que ocultan su presencia al SO.

\subsection{Capacidades del Driver}

\begin{itemize}[leftmargin=2em]
    \item \textbf{Intercepcion de procesos:} Usa \texttt{PsSetCreateProcessNotifyRoutineEx} para capturar cualquier proceso que se crea o destruye en el sistema, incluso los que se ocultan del Administrador de Tareas.
    \item \textbf{Monitoreo de imagenes:} Usa \texttt{PsSetLoadImageNotifyRoutine} para detectar la carga de DLLs y ejecutables, especialmente desde rutas inusuales.
    \item \textbf{Intercepcion de red (nivel kernel):} Puede capturar paquetes antes de que lleguen al stack de red de usuario, detectando C2 \textit{beacons} ocultos.
\end{itemize}

\subsection{Comunicacion con el Agente}

\begin{lstlisting}[language=C, caption={Fragmento del driver: comunicacion via DeviceIoControl}]
// El driver expone un dispositivo al espacio de usuario
UNICODE_STRING deviceName = RTL_CONSTANT_STRING(L"\\Device\\SatriDriver");
IoCreateDevice(DriverObject, 0, &deviceName,
               FILE_DEVICE_UNKNOWN, 0, FALSE, &gDeviceObject);

// El Kernel Bridge en Python lo abre y lee eventos
HANDLE hDevice = CreateFile("\\\\.\\SatriDriver",
    GENERIC_READ | GENERIC_WRITE, 0, NULL,
    OPEN_EXISTING, FILE_ATTRIBUTE_NORMAL, NULL);
\end{lstlisting}

\begin{notabox}
En entornos de prueba, el driver requiere activar el modo de prueba de Windows: \texttt{bcdedit /set testsigning on}. En produccion se necesita un certificado de firma de codigo EV valido.
\end{notabox}

\imagespace{Dashboard SOC: seccion de Agentes Kernel con driver Ring-0 activo}{Dashboard SOC: seccion de Agentes Kernel con driver Ring-0 activo}

\section{Extension de Navegador Chrome/Edge}

\subsection{Descripcion}
La extension de navegador es una \textbf{extension Manifest v3} para Chrome y Edge que monitorea la navegacion web en tiempo real. Su funcion principal es bloquear el acceso a sitios web cuyo dominio o IP de destino aparezca en la lista de IoCs maliciosos sincronizada desde el backend.

\subsection{Componentes de la Extension}

\begin{table}[H]
\centering
\caption{Componentes de la extension de navegador SATRI}
\begin{tabular}{@{}p{4cm}p{8cm}@{}}
\toprule
\textbf{Archivo} & \textbf{Funcion} \\
\midrule
\texttt{manifest.json}         & Declaracion de permisos, service worker \\
\texttt{background.js}         & Service Worker: sincroniza IoCs con el agente local, verifica URLs \\
\texttt{content.js}            & Script inyectado: detecta formularios de phishing y redirecciones \\
\texttt{content\_blocker.js}   & Inyecta el overlay de bloqueo en paginas maliciosas \\
\texttt{popup.html/js}         & Interfaz del popup: estado de proteccion activa \\
\bottomrule
\end{tabular}
\end{table}

\imagespace{Popup de la extension de navegador SATRI mostrando el estado de proteccion activo}{Extension de navegador SATRI: popup de estado y proteccion activa en Chrome}

% =============================================================================
\chapter{Infraestructura y Despliegue}
% =============================================================================

\section{Orquestacion con Docker Compose}

El sistema SATRI se despliega completamente con un unico archivo \texttt{docker-compose.yml} ubicado en la raiz del repositorio. Define los 14 servicios, sus dependencias de arranque (\texttt{depends\_on}), variables de entorno, puertos, volumenes persistentes y la red \texttt{satri-net}.

\begin{lstlisting}[language=bash, caption={Comandos de despliegue del entorno SATRI completo}]
# 1. Construir imagenes y levantar todos los contenedores
docker compose up -d --build

# 2. Inicializar topicos Kafka (script de configuracion)
powershell -ExecutionPolicy Bypass -File .\init.ps1

# 3. Verificar estado de todos los servicios
powershell -ExecutionPolicy Bypass -File .\health_check.ps1
\end{lstlisting}

\imagespace{Salida del script run\_all.ps1 con el inicio exitoso de todos los contenedores}{Inicio completo del entorno SATRI usando run\_all.ps1}

\section{Traefik como API Gateway}

Traefik v3 actua como punto de entrada unico y reverse proxy. Sus rutas se configuran mediante labels de Docker Compose:

\begin{table}[H]
\centering
\caption{Rutas configuradas en Traefik API Gateway}
\begin{tabular}{@{}p{3cm}p{4cm}p{5cm}@{}}
\toprule
\textbf{Prefijo / Puerto} & \textbf{Servicio Destino} & \textbf{Descripcion} \\
\midrule
\texttt{/api}       & satri-mml (8000)        & Endpoint de ingesta y gestion de logs \\
\texttt{/api/mia-vt} & satri-mia-vt (8001)    & API de inteligencia VirusTotal \\
\texttt{/soc}       & satri-dashboard (8005)  & Dashboard SOC (frontend + API) \\
\texttt{:8080}      & Traefik Dashboard       & Panel de administracion de Traefik \\
\bottomrule
\end{tabular}
\end{table}

\imagespace{Panel de administracion de Traefik (localhost:8080) mostrando los routers y servicios activos}{Traefik API Gateway: panel de administracion con rutas y servicios del sistema SATRI}

\section{Despliegue en Azure Kubernetes Service (AKS)}

Para garantizar la alta disponibilidad, resiliencia y escalabilidad del entorno, el sistema SATRI ha sido desplegado en la nube de Microsoft Azure utilizando \textbf{Azure Kubernetes Service (AKS)}. Este enfoque permite operar la plataforma bajo condiciones reales de produccion.

\subsection{Provisionamiento de Maquinas Virtuales (Nodos)}

El cluster fue provisionado con un \textit{Node Pool} conformado por \textbf{dos maquinas virtuales (nodos)}. Estas maquinas virtuales proveen el computo subyacente y permiten distribuir la carga de trabajo de los distintos componentes del sistema para evitar puntos unicos de fallo.

\imagespace{Captura del portal de Azure o terminal mostrando las 2 maquinas virtuales/nodos del cluster AKS}{Nodos del cluster AKS: provisionamiento de dos maquinas virtuales en Azure}

\subsection{Despliegue y Separacion de Microservicios}

Bajo este esquema de orquestacion, los microservicios fueron separados en \textit{Deployments} independientes de Kubernetes. Esto permitio que los \textit{Pods} (instancias de los microservicios) se alojaran y distribuyeran dinamicamente a traves de las dos maquinas virtuales disponibles. Por ejemplo, los motores de analisis y el bus de eventos (Kafka) pueden residir en nodos distintos para balancear el consumo de CPU y memoria.

\imagespace{Captura de kubectl o Azure dashboard mostrando los Pods distribuidos en diferentes Nodos (Nodes)}{Separacion de microservicios: Pods en ejecucion distribuidos a traves de las dos maquinas virtuales}

\subsection{Comunicacion en el Cluster (Red Interna)}

La comunicacion entre los microservicios, previamente manejada por la red de Docker Compose, fue implementada mediante recursos \textit{Service} de tipo \texttt{ClusterIP}. Esto dota a las maquinas virtuales de conectividad fluida para sus cargas de trabajo:
\begin{itemize}[leftmargin=2em]
    \item \textbf{Descubrimiento de Servicios:} Los microservicios se comunican entre si mediante nombres DNS internos (gestionados por CoreDNS), sin importar en que maquina virtual se encuentre fisicamente el Pod de destino.
    \item \textbf{Balanceo de Carga Interno:} El trafico se enruta de nodo a nodo automaticamente hacia los Pods saludables. Esto permite que servicios criticos como Kafka, los motores de correlacion (MCOR) y las bases de datos interactuen de forma transparente.
\end{itemize}

Esta arquitectura valida que las maquinas virtuales se comunican exitosamente a nivel de red virtual de Azure (VNet), manteniendo la cohesion y el funcionamiento del sistema SATRI en un entorno totalmente distribuido.

\imagespace{Captura de los Services (ClusterIP) o logs mostrando la comunicacion de red exitosa entre microservicios alojados en distintas VMs}{Comunicacion de red: conectividad exitosa entre microservicios a traves de los nodos del cluster}

% =============================================================================
\chapter{Modelo de Privacidad y Consentimiento}
\label{privacidad}
% =============================================================================

\section{Principio Fundamental --- Privacy by Design}

El sistema SATRI fue disenado bajo el principio de \textbf{Privacy by Design} (RFC 8280). El documento \texttt{PRIVACY\_MODEL.md} actua como un bloqueo: ningun codigo de recoleccion de datos puede escribirse hasta que el modelo sea formalmente aprobado. El marco legal de referencia es la \textbf{Ley de Proteccion de Datos Personales N\textdegree{} 29733} de Peru.

\begin{quote}
\textit{``Solo se envia lo necesario para que el panel central entienda `hubo una amenaza de tipo X con severidad Y', nunca `que estaba haciendo exactamente la persona o que contenido visito'.''} --- PRIVACY\_MODEL.md
\end{quote}

\section{Esquema del Evento Agregado}

El unico formato que el agente desktop puede enviar al backend es el siguiente esquema JSON:

\begin{lstlisting}[language=Python, caption={Esquema del evento anonimizado enviado por el agente SATRI}]
{
  "schema_version": "1.1",
  "event_id":       "evt-a1b2c3d4",
  "device_id":      "dev-hashed-uuid",          // UUID hasheado, no el real
  "timestamp":      "2026-06-18T06:00:00Z",
  "event_type":     "suspicious_process_detected",
  "severity":       "HIGH",
  "category":       "process_monitor",
  "details": {
    "resource_hash": "sha256:e3b0c4429...",     // Hash del archivo, NO el archivo
    "action_taken":  "blocked",
    "threat_name":   "Trojan.GenericKD",
    "vt_score":      "12/72",
    "geo_country":   "RU",
    "geo_city":      "Moscow"                   // Ciudad aproximada, NO GPS
  },
  "behavior_metrics": {
    "login_hour":               6,
    "day_of_week":              4,
    "session_duration_minutes": 120,
    "actions_per_minute_avg":   2.5
  }
}
\end{lstlisting}

\section{Campos Permitidos y Prohibidos}

\begin{table}[H]
\centering
\caption{Politica de datos del agente SATRI: campos permitidos y prohibidos}
\begin{tabular}{@{}p{1.5cm}p{4cm}p{6.5cm}@{}}
\toprule
\textbf{Estado} & \textbf{Campo} & \textbf{Razon} \\
\midrule
\textcolor{satriGreen}{\textbf{OK}} & \texttt{event\_type}           & Tipo de amenaza detectada \\
\textcolor{satriGreen}{\textbf{OK}} & \texttt{severity}              & Nivel de riesgo \\
\textcolor{satriGreen}{\textbf{OK}} & \texttt{details.resource\_hash} & Hash SHA-256 (no el archivo en si) \\
\textcolor{satriGreen}{\textbf{OK}} & \texttt{details.action\_taken} & Accion automatica tomada \\
\textcolor{satriGreen}{\textbf{OK}} & \texttt{behavior\_metrics.*}   & Estadisticas anonimas de sesion \\
\textcolor{satriGreen}{\textbf{OK}} & \texttt{geo\_country}          & Pais aproximado por IP (ISO) \\
\midrule
\textcolor{satriRed}{\textbf{NO}}   & \texttt{url / visited\_url}     & Revela actividad de navegacion \\
\textcolor{satriRed}{\textbf{NO}}   & \texttt{page\_content}          & Revela contenido visitado \\
\textcolor{satriRed}{\textbf{NO}}   & \texttt{window\_title}          & Revela aplicacion/documento en uso \\
\textcolor{satriRed}{\textbf{NO}}   & \texttt{keystrokes / keylog}    & Keylogging explicitamente prohibido \\
\textcolor{satriRed}{\textbf{NO}}   & \texttt{screenshot}             & Captura visual de actividad \\
\textcolor{satriRed}{\textbf{NO}}   & \texttt{file\_path / file\_name} & Revela estructura de carpetas personal \\
\textcolor{satriRed}{\textbf{NO}}   & \texttt{clipboard}              & Contenido del portapapeles \\
\textcolor{satriRed}{\textbf{NO}}   & \texttt{password / credentials} & Credenciales del usuario \\
\textcolor{satriRed}{\textbf{NO}}   & \texttt{username / hostname}    & Identificadores directos del usuario \\
\textcolor{satriRed}{\textbf{NO}}   & \texttt{gps\_location}          & Ubicacion GPS precisa \\
\bottomrule
\end{tabular}
\end{table}

\imagespace{Pantalla de onboarding del agente con el texto de consentimiento informado al usuario}{Agente SATRI: pantalla de consentimiento informado previo al inicio del monitoreo}

% =============================================================================
\chapter{Integracion Continua con Jenkins}
% =============================================================================

\section{Descripcion del Pipeline}

El proyecto implementa un pipeline de \textbf{Integracion Continua (CI)} mediante \textbf{Jenkins}, con el archivo \texttt{Jenkinsfile} en la raiz del repositorio. El pipeline se activa automaticamente en cada \textit{push} al repositorio y garantiza que ningun cambio rompa el contrato de la API antes de ser integrado.

\begin{alertbox}
El pipeline usa \textbf{Docker-out-of-Docker (DooD)} montando el socket \texttt{/var/run/docker.sock} en el contenedor del agente Jenkins. Esto es necesario para que Testcontainers pueda arrancar contenedores MongoDB efimeros durante las pruebas.
\end{alertbox}

\section{Etapas del Pipeline}

\begin{table}[H]
\centering
\caption{Etapas del pipeline Jenkins de SATRI}
\begin{tabular}{@{}p{0.5cm}p{4cm}p{8cm}@{}}
\toprule
\textbf{\#} & \textbf{Etapa} & \textbf{Descripcion} \\
\midrule
1 & Checkout         & Jenkins obtiene el codigo fuente del repositorio Git via \texttt{checkout scm} \\
2 & Install Deps     & Instala dependencias Python del microservicio \texttt{dashboard-soc} con \texttt{pip} \\
3 & Automated Tests  & Ejecuta \texttt{pytest tests/ -v} con MongoDB efimero via Testcontainers \\
\midrule
  & Post: Success    & \textit{``Pruebas exitosas. Listo para despliegue.''} \\
  & Post: Failure    & \textit{``Las pruebas fallaron. Revisar logs de Testcontainers y pytest.''} \\
  & Post: Always     & Limpieza del entorno, Ryuk destruye el contenedor MongoDB efimero \\
\bottomrule
\end{tabular}
\end{table}

\section{Jenkinsfile Completo}

\begin{lstlisting}[language=bash, caption={Jenkinsfile completo del proyecto SATRI}]
pipeline {
    agent {
        docker {
            image 'python:3.11-slim'
            // DooD: montamos el socket para que Testcontainers levante MongoDB
            args '-v /var/run/docker.sock:/var/run/docker.sock'
        }
    }

    environment {
        PYTHONUNBUFFERED = '1'
    }

    stages {
        stage('Checkout') {
            steps {
                checkout scm
            }
        }

        stage('Install Dependencies') {
            steps {
                dir('backend-aks/brain/dashboard-soc') {
                    sh '''
                        python -m pip install --upgrade pip
                        pip install -r requirements.txt
                        pip install -r requirements-test.txt
                    '''
                }
            }
        }

        stage('Automated Tests (Testcontainers)') {
            steps {
                dir('backend-aks/brain/dashboard-soc') {
                    // Testcontainers descarga mongo:6.0, ejecuta las pruebas
                    // y destruye el contenedor automaticamente al finalizar
                    sh 'pytest tests/ -v'
                }
            }
        }
    }

    post {
        always {
            echo 'Limpiando entorno. Ryuk destruye el MongoDB efimero.'
        }
        success {
            echo 'Pruebas exitosas. Listo para integracion o despliegue.'
        }
        failure {
            echo 'Las pruebas fallaron. Revisar logs de Testcontainers y pytest.'
        }
    }
}
\end{lstlisting}

\section{Pruebas con Testcontainers}

Testcontainers es una libreria que levanta contenedores Docker reales (en este caso \texttt{mongo:6.0}) durante la ejecucion de las pruebas y los destruye al terminar. Esto proporciona:

\begin{itemize}[leftmargin=2em]
    \item \textbf{Aislamiento total:} No se requiere una instalacion de MongoDB en el entorno de CI.
    \item \textbf{Reproducibilidad:} El entorno de prueba es identico en cualquier maquina del equipo o agente de Jenkins.
    \item \textbf{Limpieza automatica:} El agente Ryuk de Testcontainers garantiza que el contenedor sea eliminado incluso si las pruebas fallan abruptamente.
\end{itemize}

\imagespace{Pantalla de Jenkins con el pipeline de CI: etapas Checkout, Install, Tests en color verde (exitosas)}{Jenkins CI Pipeline: todas las etapas completadas exitosamente}

\imagespace{Salida completa de pytest -v mostrando los tests de Testcontainers con resultado PASSED}{Pruebas automatizadas: resultados de pytest con Testcontainers y MongoDB efimero}

% seguridad y evidencia

% =============================================================================
\chapter{Seguridad del Sistema}
% =============================================================================

\section{Autenticacion y Autorizacion}

\begin{itemize}[leftmargin=2em]
    \item \textbf{JWT (JSON Web Tokens):} El Dashboard SOC genera tokens JWT firmados al hacer login. Todos los endpoints requieren el token en el header \texttt{Authorization: Bearer <token>}. La expiracion se controla con \texttt{ACCESS\_TOKEN\_EXPIRE\_MINUTES}.
    \item \textbf{Bearer Tokens para Agentes:} Los agentes y servicios internos usan Bearer Tokens estaticos (configurados en \texttt{.env}) para acceder al endpoint de ingesta del MML.
    \item \textbf{MongoDB autenticado:} La base de datos requiere credenciales (\texttt{satri:satri2025}), gestionadas via variables de entorno.
    \item \textbf{HTTPS via Traefik:} En produccion, Traefik gestiona certificados TLS/SSL automaticos para cifrar todo el trafico externo.
\end{itemize}

\section{Rate Limiting y Deduplicacion}

\begin{itemize}[leftmargin=2em]
    \item \textbf{SlowAPI:} El MML aplica 300 peticiones/minuto por IP para prevenir ataques de flooding.
    \item \textbf{Filtro de Bloom:} Deteccion de eventos duplicados en tiempo $O(1)$ con \texttt{pybloom-live}, sin consultar ninguna base de datos.
    \item \textbf{Cache Redis TTL:} Resultados de VirusTotal cacheados en Redis para reducir llamadas a la API externa.
\end{itemize}

\section{Observabilidad}

\begin{itemize}[leftmargin=2em]
    \item \textbf{Prometheus:} Endpoint \texttt{/metrics} en el MML con metricas de latencia, peticiones y errores.
    \item \textbf{Loguru:} Logging estructurado en todos los microservicios con niveles DEBUG/INFO/WARNING/ERROR/CRITICAL.
    \item \textbf{OpenSearch:} Todos los logs enriquecidos indexados para busqueda forense y analisis retroactivo.
\end{itemize}

% =============================================================================
\chapter{Evidencia de Funcionamiento}
% =============================================================================

Este capitulo reserva un espacio visual para cada evidencia relevante del sistema. Insertar la captura de pantalla correspondiente reemplazando el recuadro gris.

\section{Estado del Entorno Docker}

\imagespace{Salida de docker compose ps con los 14 contenedores en estado Up}{Todos los servicios del sistema SATRI corriendo en Docker}

\section{Dashboard SOC --- Vista General}

\imagespace{Pantalla completa del Dashboard SOC con KPIs y feed de alertas}{Dashboard SOC: vista general con metricas de seguridad en tiempo real}

\section{Feed de Alertas en Tiempo Real (WebSocket)}

\imagespace{Feed lateral del Dashboard con alertas llegando instantaneamente}{Dashboard SOC: alertas en tiempo real via WebSocket sin recargar la pagina}

\section{Tabla de Incidentes y Triage}

\imagespace{Tabla de alertas con columnas: Severidad, IP, Regla MCOR, Prioridad, Accion, Botones}{Dashboard SOC: tabla de incidentes con triage automatizado y opciones de respuesta}

\section{IoC Intelligence}

\imagespace{Seccion IoC Intelligence con IPs maliciosas y su score de VirusTotal}{Dashboard SOC: Inteligencia de IoC con scores de VirusTotal}

\section{Heatmap Global de Amenazas}

\imagespace{Mapa mundial oscuro con puntos de calor sobre paises con ataques detectados}{Dashboard SOC: Heatmap Global con distribucion geografica de amenazas}

\section{Terminal de Raw Logs}

\imagespace{Terminal de logs coloreados por tipo y severidad en tiempo real}{Dashboard SOC: terminal Raw Logs con coloreado semantico por tipo de evento}

\section{Aislamiento Remoto de Host}

\imagespace{Modal de confirmacion de aislamiento de host en el Dashboard}{Dashboard SOC: aislamiento remoto de endpoint comprometido}

\section{Descarga del Agente de Proteccion}

\imagespace{Seccion de descarga de SATRI Protection con instrucciones de instalacion}{Dashboard SOC: seccion de descarga del agente de proteccion}

\section{Agente Desktop en Windows}

\imagespace{SATRI Protection.exe corriendo en el Administrador de Tareas de Windows}{Agente Desktop SATRI corriendo en segundo plano en Windows}

\section{Extension de Navegador}

\imagespace{Popup de la extension SATRI en Chrome con estado Proteccion Activa}{Extension de navegador SATRI Protection activa en Google Chrome}

\section{Jenkins Pipeline Exitoso}

\imagespace{Pantalla Jenkins: etapas Checkout, Install Dependencies y Tests en verde}{Jenkins CI: pipeline completado con todas las etapas exitosas}

\section{Pruebas con Testcontainers}

\imagespace{Salida de pytest -v mostrando todos los tests PASSED con Testcontainers}{Pytest + Testcontainers: pruebas automatizadas con MongoDB efimero exitosas}

\section{Inicio del Sistema con run\_all.ps1}

\imagespace{Consola PowerShell mostrando el inicio de todos los microservicios SATRI}{Arranque completo del entorno SATRI con run\_all.ps1}

% =============================================================================
\chapter{Estructura del Repositorio}
% =============================================================================

\begin{lstlisting}[language=bash, caption={Arbol de directorios del proyecto SATRI}]
Proyecto_Software/
|-- Jenkinsfile                     # Pipeline CI con Jenkins
|-- docker-compose.yml              # 14 servicios orquestados
|-- MCOR_REGLAS.md                  # 14 reglas MITRE ATT&CK especificadas
|-- PRIVACY_MODEL.md                # Modelo de privacidad y consentimiento
|-- README.md                       # Inicio rapido
|-- run_all.ps1                     # Script de arranque del entorno
|-- health_check.ps1                # Verificacion de estado de servicios
|-- init.ps1                        # Inicializacion de topicos Kafka
|-- .env / .env.example             # Variables de entorno y secretos
|
|-- backend-aks/
|   |-- frontline/
|   |   |-- mml/                    # MML: Ingesta, Normalizacion, GeoIP
|   |   |   |-- main.py
|   |   |   |-- normalizer.py
|   |   |   |-- geoip_enrich.py
|   |   |   |-- bloom_dedup.py
|   |   |   |-- opensearch_client.py
|   |   |   \-- Dockerfile
|   |   |
|   |   |-- mcor/                   # MCOR: Motor SIEM - 14 reglas MITRE
|   |   |   |-- main.py
|   |   |   |-- rules/
|   |   |   |   |-- regla_01_fuerza_bruta_ssh.py
|   |   |   |   |-- regla_02_escalada_sudo.py
|   |   |   |   \-- ... (14 archivos)
|   |   |   \-- Dockerfile
|   |   |
|   |   \-- mia_ml/                 # MIA-ML: UEBA / IsolationForest
|   |       |-- main.py
|   |       |-- model.py
|   |       \-- Dockerfile
|   |
|   \-- brain/
|       |-- mia_vt/                 # MIA-VT: Inteligencia VirusTotal
|       |-- mta/                    # MTA: Triage automatizado
|       |-- integrations/           # SOAR + Jira
|       \-- dashboard-soc/          # Dashboard SOC (FastAPI + SPA)
|           |-- main.py
|           |-- auth.py             # JWT generacion y verificacion
|           |-- tests/
|           |   \-- test_api.py     # Pruebas con Testcontainers
|           |-- static/
|           |   \-- index.html      # SPA (HTML + CSS + JS puro)
|           |-- requirements.txt
|           |-- requirements-test.txt
|           \-- Dockerfile
|
|-- desktop-agent/                  # Agente Desktop Python compilado
|   |-- main.py
|   |-- modules/
|   |   |-- process_monitor.py
|   |   |-- ransomware_guard.py
|   |   |-- download_scanner.py
|   |   |-- wifi_check.py
|   |   |-- network_monitor.py
|   |   |-- activity_monitor.py
|   |   |-- browser_bridge.py
|   |   \-- kernel_bridge.py
|   |-- reporter.py
|   |-- consent.py
|   |-- ioc_sync.py
|   \-- local_api.py
|
|-- browser-extension/              # Extension Chrome/Edge Manifest v3
|   |-- manifest.json
|   |-- background.js
|   |-- content.js
|   |-- content_blocker.js
|   \-- popup.html / popup.js
|
|-- agent_kernel/                   # Driver Kernel Windows (Ring-0 / C)
|   |-- satri_driver.c
|   |-- satri_driver.inf
|   |-- usermode_client.cpp
|   |-- install_driver.ps1
|   \-- README.md
|
|-- infra/
|   \-- vector/
|       \-- vector.toml             # Configuracion del recolector Vector.dev
|
\-- docs/
    \-- documentacion_satri.tex     # Este documento
\end{lstlisting}

% =============================================================================
\chapter{Conclusiones y Trabajo Futuro}
% =============================================================================

\section{Logros del Proyecto}

El sistema SATRI demuestra la implementacion exitosa de una plataforma de ciberseguridad integral de nivel empresarial:

\begin{enumerate}[leftmargin=2em]
    \item \textbf{Arquitectura Moderna y Escalable:} La implementacion de una arquitectura orientada a eventos con Apache Kafka garantiza desacoplamiento, alta disponibilidad y escalabilidad horizontal de cada componente de forma independiente.

    \item \textbf{Deteccion Multi-capa:} La combinacion de reglas deterministicas MITRE ATT\&CK (MCOR), inteligencia de amenazas con VirusTotal (MIA-VT) y Machine Learning no supervisado (MIA-ML/UEBA) proporciona cobertura de deteccion que ninguno de estos enfoques puede lograr individualmente.

    \item \textbf{Telemetria Completa:} La triada de agente de escritorio + extension de navegador + driver de kernel cubre todos los vectores de ataque modernos: procesos maliciosos, ransomware, exfiltracion de red, phishing web y amenazas de nivel kernel (rootkits).

    \item \textbf{Privacy by Design desde el Dia 1:} La aprobacion del modelo de privacidad como bloqueante antes de escribir cualquier codigo de recoleccion de datos garantiza el cumplimiento de la Ley N\textdegree{} 29733 sin comprometer la efectividad de la deteccion.

    \item \textbf{Observabilidad Total:} La integracion de Prometheus, Loguru, OpenSearch y el Dashboard SOC con WebSockets en tiempo real proporciona visibilidad completa del estado del sistema para los analistas SOC.

    \item \textbf{Calidad Garantizada por CI:} El pipeline Jenkins + Testcontainers valida automaticamente que la API del Dashboard SOC cumpla su contrato en cada cambio de codigo, con un entorno de pruebas 100\% reproducible.

    \item \textbf{Despliegue Simplificado:} La orquestacion completa con un unico \texttt{docker-compose.yml} permite desplegar los 14 servicios del sistema con un solo comando.
\end{enumerate}

\section{Trabajo Futuro}

\begin{itemize}[leftmargin=2em]
    \item \textbf{Kubernetes (AKS/EKS):} Migrar el despliegue de Docker Compose a un cluster Kubernetes gestionado para alta disponibilidad real en produccion.
    \item \textbf{Firma del Driver:} Obtener un certificado EV de firma de codigo para distribuir el driver de kernel en produccion sin requerir el modo de prueba de Windows.
    \item \textbf{RBAC en el Dashboard:} Implementar control de acceso basado en roles (RBAC) para diferenciar entre analistas SOC de nivel L1, L2 y administradores.
    \item \textbf{Integracion SIEM externa:} Exportar alertas hacia plataformas SIEM externas (Splunk, IBM QRadar) via APIs estandar (STIX/TAXII, CEF).
    \item \textbf{UEBA con datos reales:} Entrenar el modelo Isolation Forest con conjuntos de datos historicos reales de multiples organizaciones para mejorar la precision de deteccion.
    \item \textbf{HTTPS en produccion:} Configurar Traefik con Let's Encrypt para dominio HTTPS propio y despliegue en nube publica.
    \item \textbf{Soporte macOS/Linux del agente:} Extender el agente desktop a sistemas macOS y distribuciones Linux con deteccion de amenazas especificas de cada plataforma.
\end{itemize}

% =============

\end{document}